\documentclass[conference=acm]{confpaper}

% General
% -------
\usepackage{pdflscape}

% Bibliography
% ------------
\bibliography{paper}

% Commands
% --------

\pgfsetarrows{latex-latex}
\tikzset{%
  base/.style = {inner sep=5pt,
                 text centered,
                 thin,
                 font=\rmfamily},
  round/.style = {base,
                  rectangle,
                  rounded corners=1ex,
                  draw=black,
                  fill=gray!20,
                  minimum height=0.35in}
}

%%%%%%%%%%%%%%%%%%%%  SPACE SAVINGS %%%%%%%%%%%%%%%%%%%%%%

\usepackage[all=normal,wordspacing]{savetrees}

%%%%%%%%%%%%%%%%%%%%  TITLE/AUTHORS  %%%%%%%%%%%%%%%%%%%%%

\title{Bias, Omission, and Control: How Censorship Shapes LLM Responses}

\author{ 
\begin{minipage}[t]{0.45\textwidth}
\centering
Nakyung (Allison) Lee \\
nklee \\
University of Chicago \\
Chicago, IL, USA
\end{minipage}%
\hfill
\begin{minipage}[t]{0.45\textwidth}
\centering
Clarice Kim \\
claricek\\
University of Chicago \\
Chicago, IL, USA
\end{minipage}
}


\titlespacing\section{0pt}{*1.8}{*1.1}
\titlespacing\subsection{0pt}{*1.5}{*1.1}
\titlespacing\subsection{0pt}{*1.3}{*0.8}

%%%%%%%%%%%%%%%%%%%%  START OF DOCUMENT  %%%%%%%%%%%%%%%%%

\begin{document}

\abstract{
Large language models (LLMs) are increasingly being integrated into everyday
technologies, and the integrity and diversity of the data they are trained on
have become critical public concerns. Online information is not uniformly
accessible as governments, organizations, and ISPs employ content filtering,
website blocking, and other censorship methods to restrict access to specific
online content. Consequences of censorship on society and humans are generally
well documented. However, its effects on LLMs that rely on large-scale internet
data for training remain less understood.

Unlike human-facing censorship, where blocked websites or restricted content signal that information is being withheld, LLM censorship operates silently and leaves no trace. An LLM that never encountered certain information during training, cannot tell us what it doesn't know, and users also have no way to detect the absence. This creates a problem where we are building AI systems that don't just reflect censorship but naturalize it. It is important to understand how
censorship shapes what LLMs can “see” and “learn” to ensure fairness and
robustness of these models, in addition to preserving the free flow of
knowledge in artificial intelligence (AI) systems.
}

\maketitle

\section{Introduction}
This project investigates how internet censorship techniques: blocking, firewalls, flooding, and disinformation, affect the training data and outputs of LLMs. Censorship limits the model's exposure to certain kinds of online content, which could result in biased or incomplete outputs. For instance, if the environment gives higher rewards to actions that reflect social or institutional biases, the RL agent will treat those actions as optimal. This is dangerous because it can reinforce existing discrimination or social marginalization, amplifying unfair biases or discriminatory behaviors in society.

We will aim to characterize how forms of filtering, specifically firewalls and disinformation, may alter the composition of LLM training. We will then analyze how altered data inputs influence model outputs in relation to information omission, bias, and factual consistency. We will also conduct an experiment to determine whether LLMs, specifically ChatGPT-4, reflects disinformation bias using the LIAR2 Dataset. The LIAR2 Dataset is a collection of statements from internet sources (especially social media) that are labeled from 0 to 5 to represent statements that are least to most true. Our experiment will be asking ChatGPT-4 whether the statements in the LIAR2 Dataset are True or False to examine whether the LLM has been trained on biased data or disinformation circulating on the internet and how that may impact LLM outputs.

\section{Related Work}

Ahmed and Knockel \cite{impactononline} have examined how online censorship shapes large language models (2024) and Wu et al. \cite{firewall-encrypted} studied the infrastructural and technical mechanisms of internet censorship, specifically how China's Great Firewall detects and blocks fully encrypted traffic. Our project will extend these two lines of inquiry: bridging their perspectives and findings, and investigating the resulting implications of model training and data curation.
\textcolor{red}{confirm the below paragraph for correct citations and literature used}
We will review prior literature across several themes to ground our research. To examine the general impact of online censorship on LLMs, we will reference Ahmed and Knockel \cite{impactononline} and Glukhov et al. \cite{unknown2023censorshipLLMsMLOrSecurity}, who show how censorship shapes model outputs and information flow. To understand how censorship restricts data access, we will draw on Hoang et al. \cite{hoang2024howcensorshipreducesavailability} and studies on firewalls and WAFs Wu et al. \cite{firewall-encrypted}, which document technical limits on web content and resulting data gaps. To explore how censored or biased data influences LLM behavior, we will examine studies on bias and omission (Yang and Roberts, Hovy and Prabhumoye, Noels et al.) \cite{yang2021censoredbiaseddataLLMs} \cite{hovy2022fiveSourcesOfBiasNLP} \cite{noels2025whatLLMsDoNotTalkAbout} and on disinformation and data poisoning (Jiang et al., Danet, Maus et al.) \cite{jiang2023forcinggenerativemodelsdegenerate} \cite{unknown2023floodingLLMs} \cite{maus2023blackboxadversarialprompting}. Finally, to contextualize applications of censorship-like techniques for IP protection, we will review research on AI and copyright (Wei et al., Barbera) \cite{wei2024copyrighttakedownLLMs} \cite{edpb2025aiPrivacyRisksLLMs} to understand existing methods and their limitations.


\section{Disinformation}
Disinformation refers to the intentional spread of false or misleading information. In the context of large language models (LLMs), disinformation manifests through two interconnected mechanisms: data poisoning, where adversaries intentionally corrupt training data to manipulate model behavior, and disinformation reproduction, in which models unintentionally absorb and amplify misinformation present in their training corpora. Together, these mechanisms illustrate how both intentional attacks and passive exposure can undermine the reliability and trustworthiness of LLM outputs.

\subsection{Data Poisoning as a Method to Spread Disinformation}
Data poisoning represents a deliberate form of cyberattack in which adversaries intentionally corrupt an LLM's training data to manipulate its behavior in specific ways. By injecting malicious or misleading examples into the training process, adversaries can induce systematic errors, degrade model performance, or implant hidden vulnerabilities that can be exploited after deployment. As Goldblum et al. (2021) document, reinforcement learning (RL) algorithms are particularly susceptible to poisoning attacks due to their reliance on reward signals that can be strategically manipulated. RL systems can be manipulated through the alteration of reward signals, especially in contextual bandits settings where the model learns from feedback on its actions. By corrupting the reward function, adversaries can steer the model to learn incorrect associations between inputs and desired outputs. Similarly, online learning algorithms—which continuously update their parameters based on new data—are vulnerable to systematic and gradual poisoning, allowing attackers to gradually shift model behavior in desired directions without triggering detection mechanisms. \CITE{Goldblum et al., 2021}
Empirical studies further demonstrate the effectiveness and efficiency of these attacks. Liu and Shroff (2019) evaluate data poisoning strategies against stochastic multi-armed bandit algorithms, showing that attacks against $\epsilon$-greedy, Upper Confidence Bound (UCB), and Thompson Sampling methods were consistently successful. In particular, their Adaptive attack by Constant Estimation (ACE) forced the algorithm to pull the target arm linearly in time, while the attack costs remained small relative to the overall performance loss suffered by the system. %{CITE liu}
Beyond bandit and RL settings, data poisoning also poses a serious threat to recommender systems. Fang et al. (2020) demonstrate that data poisoning attacks can successfully promote a specific target item. They found that injecting as little as 0.5% fake users into the Yelp dataset enabled their constructed poisoning attacks to promote a randomly selected target item into the top-N recommendation list for 150 times more normal users. \CITE{Fang} These findings underscore how minimal adversarial effort can yield outsized influence over information exposure at scale.
Recent work extends these concerns to natural language generation (NLG) tasks directly relevant to LLMs. Jiang et al. (2023) investigate data poisoning attacks on text summarization and text completion, finding that under “fixed” trigger insertion and full fine-tuning, poisoning only 1% of training data was sufficient to successfully compromise summarization models, while text completion required only 5% of poisoned data. \CITE These results highlight how even small manipulations in training corpora can significantly alter model behavior in high-stakes language tasks.
Taken together, data poisoning in the context of training-only attacks can reshape LLM outputs through several interrelated pathways, each with significant implications for model reliability and safety. First, degraded performance occurs as poisoned data interferes with the model’s ability to learn accurate patterns from legitimate training examples. Second, misclassification emerges when models systematically produce incorrect but plausible outputs. Third, backdoors can be embedded through hidden triggers in the model that cause malicious behavior when activated. These backdoors may remain dormant during normal operation but can be exploited through specific input patterns known only to the attacker. Finally, bias amplification can result when the model internalizes skewed patterns that reinforce harmful stereotypes or discriminatory outcomes. Particularly troubling is that these effects may remain undetected during surface-level evaluations, while still perpetuating social harms. When combined with broader information controls or censorship mechanisms, these impacts can compound in ways that are difficult to detect and remediate (Goldblum et al., 2021).


\section{Evaluation}
For our evaluation of censorship techniques and their impact on LLMs, we will analyze and then compare each of the different techniques. To guide our evaluation, we have generated questions such as: Are the outputs of LLMs different for each mode of censorship? Is there a certain threshold of input data bias the model can filter? Do the different censorship techniques change the way each model processes the data? While quantitative answers may not exist for some of these questions, comparing different modes of censorship will help provide reference points for our evaluation. We plan to make our comparisons more concrete by researching measurements for the degree of censorship on input data. We can similarly research ways to measure the outputs of LLMs.

As for our experimental study, we will evaluate the results in a quantitative manner by categorizing similar types of behavior in model responses and counting the occurrence within each category. Moreover, we can quantify the degree of censorship applied to our inputs by calculating the percentage of change in training data between different trials. We can qualitatively evaluate our study by analyzing whether it matches previous research—in other words, if it is rational to predict a specific model output given the variables we changed and kept the same. 

Finally, we can briefly explore the implications of our study on protecting intellectual property, particularly seeing if specific censorship methods we cited in our research can be repurposed to protect certain online content from being outputted by an LLM. 


\section{Ethics}
Regarding the ethical implications of our project, our models may repeat the same types of bias that exists in censored information because we are testing how censorship may impact LLMs. As an example, they may leave out certain opinions or topics that censorship typically hides. This is especially important for our experiment, where we expose the model to examples of censored information—not actually train it on real-life censored training data. Because this exposure does not fully simulate how a model learns from censored training data, it introduces a limitation: the effects we observe may reflect short-term bias from exposure rather than deeper, structural bias that would appear through full training. This means our results could be partial or misleading if interpreted without caution. We also risk reinforcing the same biases we are trying to study, since the model’s outputs may already be influenced by the censored inputs we provide. To address this, we will clearly document our process, interpret results carefully, and emphasize that our findings are exploratory rather than definitive.

\label{lastpage}

\microtypesetup{protrusion=false}

\newpage
\printbibliography

\vfill
\pagebreak

\end{document}

%%%%%%%%%%%%%%%%%%%%  END OF DOCUMENT  %%%%%%%%%%%%%%%%%%%%
